\input{../tex-inputs/latex-korrekturansicht-vorspann}

\section[Arthur Schnitzler an Berta Zuckerkandl, 18. 4. 1911]{L03981 Arthur Schnitzler an Berta Zuckerkandl, 18. 4. 1911}
\nopagebreak\mylabel{L03981v}
\rehead{ }\normalsize\beginnumbering\briefempfaengerindex{Zuckerkandl, Berta@\textsc{Zuckerkandl, Berta}!zzzSchnitzler, Arthur@\emph{von Arthur Schnitzler}!1911-04-181@{18. 4. 1911}|(be}
\toendnotes[C]{\smallbreak\pagebreak[2]}
\correspDesc{Versand  durch Arthur Schnitzler am 18. 4. 1911 in Menton
\newline{}Erhalt  durch Berta Zuckerkandl im Zeitraum [19. 4. 1911
                  – 23. 4. 1911?] in Wien}\toendnotes[C]{\smallbreak}
\Standort{Wien, Österreichische Nationalbibliothek, 405/B78/1 LIT MAG.}
\physDesc{Brief, 1 Blatt, 4 Seiten, 1760 Zeichen
\newline{}Handschrift: schwarze Tinte, lateinische Kurrent}\toendnotes[C]{\smallbreak}
\pstart
           \raggedleft{}{\pb}\textcolor{pink}{\textcolor{gray}{\textbf{GRAND HOTEL NATIONAL}}}\oindex{Grand Hôtel National [Menton]@\textbf{Grand Hôtel National [Menton]}, \emph{Hotel}|pw}{}\ledrightnote{\textcolor{pink}{Grand Hôtel National [Menton]}}\pend
           
\pstart
           \raggedleft{}\textcolor{pink}{\textcolor{gray}{\textbf{MENTON}}}\oindex{Menton@\textbf{Menton}|pw}{}\ledrightnote{\textcolor{pink}{Menton}}\pend
           
\pstart
           \raggedleft{}18. 4. 1911.\pend
           
\pstart{}Verehrte gnädige Frau,\pend\vspace{0.5em}
\pstart
           ich danke Ihnen sehr für Ihren Brief. \textcolor{green}{Medardus}\pwindex{Schnitzler, Arthur 15. 5. 1862 Wien – 21. 10. 1931 ebd.@\textsc{Schnitzler, Arthur} (15. 5. 1862 Wien – 21. 10. 1931 ebd.), \emph{Schriftsteller, Mediziner}!junge Medardus. Dramatische Historie in einem Vorspiel und fünf Aufzügen@\strich\emph{Der junge Medardus. Dramatische Historie in einem Vorspiel und fünf Aufzügen}|pw}{}\ledrightnote{\textcolor{green}{Der junge Medardus. Dramatische Historie in einem Vorspiel und fünf Aufzügen}}
               hab ich für \textcolor{pink}{Frankreich}\oindex{Frankreich@\textbf{Frankreich}|pw}{}\ledrightnote{\textcolor{pink}{Frankreich}} noch nicht vergeben – und
                  we{\geminationn} Sie glauben etwas damit {\kaufmannsund} dafür thun zu kö{\geminationn}en, so
               werde ich höchst einverstanden sein. Nur theil ich Ihre Hoffnungen vorläufig gar
               nicht – wobei ich \substVorne{}\textsuperscript{mit}\substDazwischen{}von\substHinten{} meinen bisherigen Erfahrungen in \textcolor{pink}{Frankreich}\oindex{Frankreich@\textbf{Frankreich}|pw}{}\ledrightnote{\textcolor{pink}{Frankreich}} ganz absehen will. Speziell aber der \textcolor{green}{Medardus}\pwindex{Schnitzler, Arthur 15. 5. 1862 Wien – 21. 10. 1931 ebd.@\textsc{Schnitzler, Arthur} (15. 5. 1862 Wien – 21. 10. 1931 ebd.), \emph{Schriftsteller, Mediziner}!junge Medardus. Dramatische Historie in einem Vorspiel und fünf Aufzügen@\strich\emph{Der junge Medardus. Dramatische Historie in einem Vorspiel und fünf Aufzügen}|pw}{}\ledrightnote{\textcolor{green}{Der junge Medardus. Dramatische Historie in einem Vorspiel und fünf Aufzügen}} – welcher \textcolor{pink}{französische}\oindex{Frankreich@\textbf{Frankreich}|pw}{}\ledrightnote{\textcolor{pink}{Frankreich}} Director wird sich dieser Mühe unterziehen? \textcolor{blue}{Antoine}\pwindex{Antoine, André 31.\,1.\,1858 Limoges – 23.\,10.\,1943 Le Pouliguen@\textsc{Antoine, André} (31.\,1.\,1858 Limoges – 23.\,10.\,1943 Le Pouliguen), \emph{Theaterleiter, Schauspieler}|pw}{}\ledrightnote{\textcolor{blue}{André Antoine}}?? Er hat von mir schon \label{K_L03981-1v}\edtext{zwei \textcolor{green}{Einacter}\pwindex{Schnitzler, Arthur 15. 5. 1862 Wien – 21. 10. 1931 ebd.@\textsc{Schnitzler, Arthur} (15. 5. 1862 Wien – 21. 10. 1931 ebd.), \emph{Schriftsteller, Mediziner}!grüne Kakadu. Groteske in einem Akt@\strich\emph{Der grüne Kakadu. Groteske in einem Akt}|pwv}\pwindex{Schnitzler, Arthur 15. 5. 1862 Wien – 21. 10. 1931 ebd.@\textsc{Schnitzler, Arthur} (15. 5. 1862 Wien – 21. 10. 1931 ebd.), \emph{Schriftsteller, Mediziner}!Gefährtin. Schauspiel in einem Akt@\strich\emph{Die Gefährtin. Schauspiel in einem Akt}|pwv}{}\ledrightnote{{$\rightarrow$}\emph{\textcolor{green}{Der grüne Kakadu. Groteske in einem Akt}}{\newline}{$\rightarrow$}\emph{\textcolor{green}{Die Gefährtin. Schauspiel in einem Akt}}} aufgeführt}{\lemma{\textnormal{\emph{zwei Einacter aufgeführt}}}\Cendnote{\textnormal{Die \textcolor{green}{Übersetzung}\pwindex{Schnitzler, Arthur 15. 5. 1862 Wien – 21. 10. 1931 ebd.@\textsc{Schnitzler, Arthur} (15. 5. 1862 Wien – 21. 10. 1931 ebd.), \emph{Schriftsteller, Mediziner}!Compagne. Comédie en une acte@\strich\emph{La Compagne. Comédie en une acte}|pwkv} von \emph{\textcolor{green}{Die Gefährtin}\pwindex{Schnitzler, Arthur 15. 5. 1862 Wien – 21. 10. 1931 ebd.@\textsc{Schnitzler, Arthur} (15. 5. 1862 Wien – 21. 10. 1931 ebd.), \emph{Schriftsteller, Mediziner}!Gefährtin. Schauspiel in einem Akt@\strich\emph{Die Gefährtin. Schauspiel in einem Akt}|pwk}} hatte
                  am 29. 4. 1902 am \emph{\textcolor{brown}{Théâtre Antoine}\orgindex{Théâtre Antoine@Théâtre Antoine|pwk}}{ }\textcolor{violet}{Premiere}\eventindex{Théâtre Antoine-Simone Berriau@\textbf{Théâtre Antoine-Simone Berriau}!Premiere von La Compagne, 29.4.1902@Premiere von La Compagne, 29.4.1902|pwkv} und wurde dort
                  insgesamt vier mal aufgeführt, die \textcolor{green}{Übersetzung}\pwindex{Schnitzler, Arthur 15. 5. 1862 Wien – 21. 10. 1931 ebd.@\textsc{Schnitzler, Arthur} (15. 5. 1862 Wien – 21. 10. 1931 ebd.), \emph{Schriftsteller, Mediziner}!Au Perroquet Vert@\strich\emph{Au Perroquet Vert}|pwkv} von \emph{\textcolor{green}{Der grüne
                     Kakadu}\pwindex{Schnitzler, Arthur 15. 5. 1862 Wien – 21. 10. 1931 ebd.@\textsc{Schnitzler, Arthur} (15. 5. 1862 Wien – 21. 10. 1931 ebd.), \emph{Schriftsteller, Mediziner}!grüne Kakadu. Groteske in einem Akt@\strich\emph{Der grüne Kakadu. Groteske in einem Akt}|pwk}} feierte am 7. 11. 1903{ }\textcolor{pink}{ebendort}\oindex{Théâtre Antoine-Simone Berriau@\textbf{Théâtre Antoine-Simone Berriau}, \emph{Theater}|pwkv}{ }\textcolor{violet}{Premiere}\eventindex{Théâtre Antoine-Simone Berriau@\textbf{Théâtre Antoine-Simone Berriau}!Premiere von Au Perroquet Vert, 7.11.1903@Premiere von Au Perroquet Vert, 7.11.1903|pwkv} und erlebte
                  insgesamt zwölf Aufführungen.}}}\label{K_L03981-1}: \textcolor{green}{Gefährtin}\pwindex{Schnitzler, Arthur 15. 5. 1862 Wien – 21. 10. 1931 ebd.@\textsc{Schnitzler, Arthur} (15. 5. 1862 Wien – 21. 10. 1931 ebd.), \emph{Schriftsteller, Mediziner}!Gefährtin. Schauspiel in einem Akt@\strich\emph{Die Gefährtin. Schauspiel in einem Akt}|pw}{}\ledrightnote{\textcolor{green}{Die Gefährtin. Schauspiel in einem Akt}}{ }{\kaufmannsund}{ }\textcolor{green}{Kakadu}\pwindex{Schnitzler, Arthur 15. 5. 1862 Wien – 21. 10. 1931 ebd.@\textsc{Schnitzler, Arthur} (15. 5. 1862 Wien – 21. 10. 1931 ebd.), \emph{Schriftsteller, Mediziner}!grüne Kakadu. Groteske in einem Akt@\strich\emph{Der grüne Kakadu. Groteske in einem Akt}|pw}{}\ledrightnote{\textcolor{green}{Der grüne Kakadu. Groteske in einem Akt}}, {\pb}mit gutem aber nicht dauernden
               Erfolge, der Cyclus \label{K_L03981-2v}\edtext{\textcolor{green}{Lebendige Stunden}\pwindex{Schnitzler, Arthur 15. 5. 1862 Wien – 21. 10. 1931 ebd.@\textsc{Schnitzler, Arthur} (15. 5. 1862 Wien – 21. 10. 1931 ebd.), \emph{Schriftsteller, Mediziner}!Lebendige Stunden. Vier Einakter@\strich\emph{Lebendige Stunden. Vier Einakter}|pw}{}\ledrightnote{\textcolor{green}{Lebendige Stunden. Vier Einakter}}}{\lemma{\textnormal{\emph{Lebendige Stunden}}}\Cendnote{\textnormal{Von diesen Übersetzungen wurden drei
                  publiziert: \emph{\textcolor{green}{La Femme au poignard}\pwindex{Schnitzler, Arthur 15. 5. 1862 Wien – 21. 10. 1931 ebd.@\textsc{Schnitzler, Arthur} (15. 5. 1862 Wien – 21. 10. 1931 ebd.), \emph{Schriftsteller, Mediziner}!femme au poignard@\strich\emph{La femme au poignard}|pwk}}. In: \emph{\textcolor{green}{Revue de Paris}\pwindex{Revue de Paris@\emph{La Revue de Paris}|pwk}}, Jg. 19, Reihe 3
                     (Mai–Juni), 15. 5. 1912, S. 225–238, \emph{\textcolor{green}{Les Derniers masques. Comédie en un act}\pwindex{Schnitzler, Arthur 15. 5. 1862 Wien – 21. 10. 1931 ebd.@\textsc{Schnitzler, Arthur} (15. 5. 1862 Wien – 21. 10. 1931 ebd.), \emph{Schriftsteller, Mediziner}!Derniers masques. Comédie en un act@\strich\emph{Les Derniers masques. Comédie en un act}|pwk}}.
                     In: \emph{\textcolor{green}{Revue Politique et Littéraire. Revue
                        bleue}\pwindex{Revue politique et littéraire@\emph{La Revue politique et littéraire}|pwk}}, Jg. 50, 2. Semester, Nr. 20, 11. 11. 1912,
                     S. 618–622; Nr. 21, 23. 11. 1912, S. 652–657 und \emph{\textcolor{green}{Littérature. Comédi en en act}\pwindex{Schnitzler, Arthur 15. 5. 1862 Wien – 21. 10. 1931 ebd.@\textsc{Schnitzler, Arthur} (15. 5. 1862 Wien – 21. 10. 1931 ebd.), \emph{Schriftsteller, Mediziner}!Littérature. Comédie en en act@\strich\emph{Littérature. Comédie en en act}|pwk}}. In: \emph{\textcolor{green}{La Revue bleue. La Revue politique et
                        littéraire}\pwindex{Revue politique et littéraire@\emph{La Revue politique et littéraire}|pwk}}, Jahrgang 52, 1. Semester, Nr. 1, 3. 1. 1914,
                     S. 11–16; Nr. 2, 10. 1. 1914, S. 44–50.}}}\label{K_L03981-2} (übersetzt
               von \textcolor{blue}{Rémon}\pwindex{Rémon, Maurice 27.\,11.\,1861 Paris – 20.\,6.\,1945 Mérignac@\textsc{Rémon, Maurice} (27.\,11.\,1861 Paris – 20.\,6.\,1945 Mérignac), \emph{Übersetzer}|pw}{}\ledrightnote{\textcolor{blue}{Maurice Rémon}} u Mme \textcolor{blue}{Valentin}\pwindex{Valentin, Noémi @\textsc{Valentin, Noémi}, \emph{Übersetzerin}|pw}{}\ledrightnote{\textcolor{blue}{Noémi Valentin}}) liegt seit etwa 6 Jahren angenommen
               bei ihm, aber er denkt nicht daran die \textcolor{green}{Sachen}\pwindex{Schnitzler, Arthur 15. 5. 1862 Wien – 21. 10. 1931 ebd.@\textsc{Schnitzler, Arthur} (15. 5. 1862 Wien – 21. 10. 1931 ebd.), \emph{Schriftsteller, Mediziner}!Lebendige Stunden. Vier Einakter@\strich\emph{Lebendige Stunden. Vier Einakter}|pwv}{}\ledrightnote{{$\rightarrow$}\emph{\textcolor{green}{Lebendige Stunden. Vier Einakter}}} aufzuführen. Nach einem Brief von \textcolor{blue}{Lugne Poë}\pwindex{Lugné-Poe, Aurélien-Marie 27.\,12.\,1869 Paris – 19.\,6.\,1940 Villeneuve-les-Avignon@\textsc{Lugné-Poe, Aurélien-Marie} (27.\,12.\,1869 Paris – 19.\,6.\,1940 Villeneuve-les-Avignon), \emph{Theaterleiter, Regisseur, Schauspieler}|pw}{}\ledrightnote{\textcolor{blue}{Aurélien-Marie Lugné-Poe}} an mich (anläßlich »\textcolor{green}{Liebelei}\pwindex{Schnitzler, Arthur 15. 5. 1862 Wien – 21. 10. 1931 ebd.@\textsc{Schnitzler, Arthur} (15. 5. 1862 Wien – 21. 10. 1931 ebd.), \emph{Schriftsteller, Mediziner}!Liebelei. Schauspiel in drei Akten@\strich\emph{Liebelei. Schauspiel in drei Akten}|pw}{}\ledrightnote{\textcolor{green}{Liebelei. Schauspiel in drei Akten}}« für die sich er und seine \textcolor{blue}{Frau}\pwindex{Desprès, Suzanne 18.\,12.\,1875 Verdun – 29.\,6.\,1951 Paris@\textsc{Desprès, Suzanne} (18.\,12.\,1875 Verdun – 29.\,6.\,1951 Paris), \emph{Schauspielerin}|pwv}{}\ledrightnote{{$\rightarrow$}\emph{\textcolor{blue}{Suzanne Desprès}}} (derer Name mir in diesem Moment absolut nicht
               einfallen will) interessirt haben, stehen jetzt die Chancen für deutsche Dichter
               recht übel in \textcolor{pink}{Frankreich}\oindex{Frankreich@\textbf{Frankreich}|pw}{}\ledrightnote{\textcolor{pink}{Frankreich}}.\pend
           
\pstart
           – Nachdem ich Ihnen auf
               diese Weise, verehrte gnädige Frau, den zu einem solchen Unternehmen nöthigen Muth
               eingeflößt habe, ka{\geminationn} ich nur wiederholen –: we{\geminationn} Sie es wagen wollen –\pend
           
\pstart
           {\pb}Jedenfalls werde ich bitten, in \textcolor{pink}{Wien}\oindex{Wien@\textbf{Wien}, \emph{Verwaltungsgebiet}|pw}{}\ledrightnote{\textcolor{pink}{Wien}} (wo ich \label{K_L03981-3v}\edtext{Anfang Mai}{\lemma{\textnormal{\emph{Anfang Mai}}}\Cendnote{\textnormal{\textcolor{blue}{Schnitzler} hatte \textcolor{pink}{Wien}\oindex{Wien@\textbf{Wien}, \emph{Verwaltungsgebiet}|pwk} am 10. 4. 1911 verlassen zu einer Reise über \textcolor{pink}{München}\oindex{München@\textbf{München}|pwk} und \textcolor{pink}{Garmisch-Partenkirchen}\oindex{Garmisch-Partenkirchen@\textbf{Garmisch-Partenkirchen}, \emph{Hauptstadt}|pwk} nach \textcolor{pink}{Ligurien}\oindex{Ligurien@\textbf{Ligurien}|pwk} und \textcolor{pink}{Menton}\oindex{Menton@\textbf{Menton}|pwk} mit Abstechern
                  nach \textcolor{pink}{Monte Carlo}\oindex{Monte Carlo@\textbf{Monte Carlo}, \emph{Ehemaliger Ort}|pwk} und \textcolor{pink}{Nizza}\oindex{Nizza@\textbf{Nizza}, \emph{Hauptstadt}|pwk}. Am 3. 5. 1911 kehrte er nach hause zurück, wo es
                  laut \emph{\textcolor{green}{Tagebuch}\pwindex{Schnitzler, Arthur 15. 5. 1862 Wien – 21. 10. 1931 ebd.@\textsc{Schnitzler, Arthur} (15. 5. 1862 Wien – 21. 10. 1931 ebd.), \emph{Schriftsteller, Mediziner}!Tagebuch@\strich\emph{Tagebuch}|pwk}} am 10. 5. 1911 zur
                  persönlichen Besprechung kam: »Bei der Hofr. \textcolor{blue}{Zuckerkandl}\pwindex{Zuckerkandl, Berta 13.\,4.\,1864 Wien – 16.\,10.\,1945 Paris@\textsc{Zuckerkandl, Berta} (13.\,4.\,1864 Wien – 16.\,10.\,1945 Paris), \emph{Schriftstellerin, Journalistin, Übersetzerin}|pw} (die mir wegen \textcolor{blue}{Antoine}\pwindex{Antoine, André 31.\,1.\,1858 Limoges – 23.\,10.\,1943 Le Pouliguen@\textsc{Antoine, André} (31.\,1.\,1858 Limoges – 23.\,10.\,1943 Le Pouliguen), \emph{Theaterleiter, Schauspieler}|pw} – \textcolor{green}{Medardus}\pwindex{Schnitzler, Arthur 15. 5. 1862 Wien – 21. 10. 1931 ebd.@\textsc{Schnitzler, Arthur} (15. 5. 1862 Wien – 21. 10. 1931 ebd.), \emph{Schriftsteller, Mediziner}!junge Medardus. Dramatische Historie in einem Vorspiel und fünf Aufzügen@\strich\emph{Der junge Medardus. Dramatische Historie in einem Vorspiel und fünf Aufzügen}|pw}
                     geschrieben). Über meine bisherigen Erfahrungen und Chancen in \textcolor{pink}{Frankreich}\oindex{Frankreich@\textbf{Frankreich}|pw}.«}}}\label{K_L03981-3} zu sein hoffe) persönlich
               über diese, und auch die \textcolor{green}{Weite
               Land}\pwindex{Schnitzler, Arthur 15. 5. 1862 Wien – 21. 10. 1931 ebd.@\textsc{Schnitzler, Arthur} (15. 5. 1862 Wien – 21. 10. 1931 ebd.), \emph{Schriftsteller, Mediziner}!weite Land. Tragikomödie in fünf Akten@\strich\emph{Das weite Land. Tragikomödie in fünf Akten}|pw}{}\ledrightnote{\textcolor{green}{Das weite Land. Tragikomödie in fünf Akten}}-Angelegenheit sprechen zu dürfen. Dieses \textcolor{green}{Stück}\pwindex{Schnitzler, Arthur 15. 5. 1862 Wien – 21. 10. 1931 ebd.@\textsc{Schnitzler, Arthur} (15. 5. 1862 Wien – 21. 10. 1931 ebd.), \emph{Schriftsteller, Mediziner}!weite Land. Tragikomödie in fünf Akten@\strich\emph{Das weite Land. Tragikomödie in fünf Akten}|pwv}{}\ledrightnote{{$\rightarrow$}\emph{\textcolor{green}{Das weite Land. Tragikomödie in fünf Akten}}} scheint mir nach den internationalen Seite mehr zu
               versprechen als der \textcolor{green}{Medardus}\pwindex{Schnitzler, Arthur 15. 5. 1862 Wien – 21. 10. 1931 ebd.@\textsc{Schnitzler, Arthur} (15. 5. 1862 Wien – 21. 10. 1931 ebd.), \emph{Schriftsteller, Mediziner}!junge Medardus. Dramatische Historie in einem Vorspiel und fünf Aufzügen@\strich\emph{Der junge Medardus. Dramatische Historie in einem Vorspiel und fünf Aufzügen}|pw}{}\ledrightnote{\textcolor{green}{Der junge Medardus. Dramatische Historie in einem Vorspiel und fünf Aufzügen}}. Man ist sowohl von
                  \textcolor{pink}{England}\oindex{England@\textbf{England}, \emph{Land}|pw}{}\ledrightnote{\textcolor{pink}{England}} als von \textcolor{pink}{Frankreich}\oindex{Frankreich@\textbf{Frankreich}|pw}{}\ledrightnote{\textcolor{pink}{Frankreich}} her (ohne \substVorne{}\textsuperscript{das Stück}\substDazwischen{}es\substHinten{} zu kennen) wegen dieses \textcolor{green}{Stücks}\pwindex{Schnitzler, Arthur 15. 5. 1862 Wien – 21. 10. 1931 ebd.@\textsc{Schnitzler, Arthur} (15. 5. 1862 Wien – 21. 10. 1931 ebd.), \emph{Schriftsteller, Mediziner}!weite Land. Tragikomödie in fünf Akten@\strich\emph{Das weite Land. Tragikomödie in fünf Akten}|pwv}{}\ledrightnote{{$\rightarrow$}\emph{\textcolor{green}{Das weite Land. Tragikomödie in fünf Akten}}} an mich herangetreten, ich habe mich aber noch nicht gebunden. Ihr
               Interesse, verehrte gnädige Frau, für meine Arbeiten ist mir in jedem Falle sehr
                  erfreu{\pb}lich; ich darf wohl fernere
               Nachrichten von Ihnen erwarten.\pend
           
\pstart
           mit wiederholten Dank und der Versicherung meiner aufrichtigen
               Hochachtung{\\[\baselineskip]}Ihr sehr ergebener{\\[\baselineskip]}\spacefill\mbox{Arthur Schnitzler}\pend
           \leftskip=0em{}\selectlanguage{ngerman}\endnumbering\briefempfaengerindex{Zuckerkandl, Berta@\textsc{Zuckerkandl, Berta}!zzzSchnitzler, Arthur@\emph{von Arthur Schnitzler}!1911-04-181@{18. 4. 1911}|)be}\mylabel{L03981h}  \newcommand{\dateiname}{L03981}\newcommand{\titel}{Arthur Schnitzler an Berta Zuckerkandl, 18. 4. 1911}\newcommand{\editorInnen}{Herausgegeben von Jahnke, SelmaMüller, Martin Anton}\input{../tex-inputs/latex-korrekturansicht-abspann}
